% Generated by scripts/gen_blueprint.py from blueprint/nodes/*.toml. Do not edit.

\bibchapter
\begin{thebibliography}{RW97}
\bibitem[Bak66]{Bak66} A.~Baker, Linear forms in the logarithms of algebraic numbers I, \emph{Mathematika} 13 (1966), 204--216.
\bibitem[Bak75]{Bak75} A.~Baker, \emph{Transcendental Number Theory}, Cambridge University Press, 1975.
\bibitem[BM80]{BM80} D.~Bertrand and D.~W.~Masser, Linear forms in elliptic integrals, \emph{Invent. Math.} 58 (1980), 283--288.
\bibitem[Bro74]{Bro74} W.~D.~Brownawell, The algebraic independence of certain numbers related by the exponential function, \emph{J. Number Theory} 6 (1974), 22--31.
\bibitem[DK24]{DK24} S.~Dasgupta and M.~Kakde, Ranks of matrices of logarithms of algebraic numbers II: the Matrix Coefficient Conjecture, arXiv:2408.08178, 2024.
\bibitem[Gel34]{Gel34} A.~O.~Gel'fond, Sur le septième problème de Hilbert, \emph{Izv. Akad. Nauk SSSR} (1934), 623--634.
\bibitem[Gel52]{Gel52} A.~O.~Gel'fond, \emph{Transcendental and Algebraic Numbers}, 1952.
\bibitem[Her1873]{Her1873} C.~Hermite, Sur la fonction exponentielle, \emph{C. R. Acad. Sci. Paris} 77 (1873), 18--24, 74--79, 226--233, 285--293.
\bibitem[KW26]{KW26} M.~Karatarakis and F.~Wiedijk, A formalization of the Gelfond--Schneider theorem, arXiv:2603.24823, 2026.
\bibitem[Lang66]{Lang66} S.~Lang, \emph{Introduction to Transcendental Numbers}, Addison-Wesley, 1966.
\bibitem[Lin1882]{Lin1882} F.~Lindemann, Über die Zahl $\pi$, \emph{Math. Ann.} 20 (1882), 213--225.
\bibitem[Mas81]{Mas81} D.~W.~Masser, A note on Baker's theorem, in \emph{Recent Progress in Analytic Number Theory}, Vol.~2 (Durham, 1979), Academic Press, 1981, 153--158.
\bibitem[Ram68]{Ram68} K.~Ramachandra, Contributions to the theory of transcendental numbers I, \emph{Acta Arith.} 14 (1968), 65--72.
\bibitem[RW95]{RW95} D.~Roy and M.~Waldschmidt, Quadratic relations between logarithms of algebraic numbers, \emph{Proc. Japan Acad. Ser. A} 71 (1995), 151--153.
\bibitem[RW97]{RW97} D.~Roy and M.~Waldschmidt, Approximation diophantienne et indépendance algébrique de logarithmes, \emph{Ann. Sci. École Norm. Sup. (4)} 30 (1997), 753--796.
\bibitem[Sch34]{Sch34} Th.~Schneider, Transzendenzuntersuchungen periodischer Funktionen, \emph{J. Reine Angew. Math.} 172 (1934), 65--69.
\bibitem[Sch57]{Sch57} Th.~Schneider, \emph{Einführung in die transzendenten Zahlen}, Springer, Berlin, 1957.
\bibitem[Tij71]{Tij71} R.~Tijdeman, \emph{Proc. Kon. Nederl. Akad. Wetensch. Ser. A} 74 (= \emph{Indag. Math.} 33) (1971), 1--7.
\bibitem[Wal71]{W71} M.~Waldschmidt, Indépendance algébrique des valeurs de la fonction exponentielle, \emph{Bull. Soc. Math. France} 99 (1971), 285--304.
\bibitem[Wal73]{W73} M.~Waldschmidt, Solution du huitième problème de Schneider, \emph{J. Number Theory} 5 (1973), 191--202.
\bibitem[Wal00]{Wal00} M.~Waldschmidt, \emph{Diophantine Approximation on Linear Algebraic Groups}, Grundlehren der mathematischen Wissenschaften 326, Springer, 2000.
\bibitem[Wal09]{Wal09} M.~Waldschmidt, Auxiliary functions in transcendence proofs, arXiv:0908.4024, 2009.
\bibitem[Wei1885]{Wei1885} K.~Weierstrass, Zu Lindemann's Abhandlung ``Über die Ludolph'sche Zahl'', \emph{Sitzungsber. Preuss. Akad. Wiss.} (1885), 1067--1085.
\end{thebibliography}
